\documentclass[11pt,a4paper]{article}
\usepackage[utf8]{inputenc}
\usepackage[spanish]{babel}
\usepackage{amsmath}
\usepackage{amsfonts}
\usepackage{amssymb}
\usepackage{geometry}
\usepackage{booktabs}
\usepackage{hyperref}
\usepackage{microtype}
\usepackage{xcolor}

\geometry{top=2.5cm, bottom=2.5cm, left=2.5cm, right=2.5cm}

\title{\textbf{Desarrollo Matemático Completo del Modelo de Segregación en Dunas con Deposición Periódica}}
\author{Influencia de la Segregación Sedimentaria en la Dinámica Hidráulica de Dunas}
\date{\today}

\begin{document}

\maketitle

\tableofcontents
\newpage

\section{Introducción}
Este documento presenta el desarrollo matemático riguroso detrás del script \texttt{duna\_bidispersa\_abanico.py}. El modelo simula el transporte advectivo y la segregación de un sedimento bidisperso (mezcla de partículas finas y gruesas) en el interior de una duna que migra a velocidad constante $c_{\text{mig}}$. El resolvedor opera en un sistema de coordenadas adaptativas sigma ($\eta$) en el marco móvil, incorporando una modulación periódica temporal y espacial del flujo en el sotavento para modelar la deposición periódica por capas (estratificación cruzada).

---

\section{Geometría de la Duna y Topografía}
La duna se define en un dominio bidimensional de longitud $L$ con condiciones de contorno periódicas en la dirección horizontal $x$. La topografía del lecho, denotada por $h(x)$, se modela mediante un perfil lineal por tramos con una cresta aguda en $x_{\text{crest}}$:

\begin{equation}
h(x) = \begin{cases} 
H_{\text{base}} + H_d \left( \frac{x}{x_{\text{crest}}} \right), & \text{si } 0 \le x \le x_{\text{crest}} \quad \text{(Barlovento / Stoss)} \\
H_{\text{base}} + H_d \left( 1 - \frac{x - x_{\text{crest}}}{L - x_{\text{crest}}} \right), & \text{si } x_{\text{crest}} < x \le L \quad \text{(Sotavento / Lee)}
\end{cases}
\end{equation}

donde $H_{\text{base}}$ es la altura del lecho base, $H_d$ es la altura neta de la duna y $x_{\text{crest}}$ es la posición de la cresta.
La derivada espacial de la topografía, $\frac{dh}{dx}$, es constante por tramos:
\begin{equation}
\frac{dh}{dx} = \begin{cases} 
\frac{H_d}{x_{\text{crest}}} > 0, & \text{si } 0 \le x \le x_{\text{crest}} \\
-\frac{H_d}{L - x_{\text{crest}}} < 0, & \text{si } x_{\text{crest}} < x \le L
\end{cases}
\end{equation}

---

\section{Transformación de Coordenadas Sigma (\(\eta\))}
Para resolver las ecuaciones en un dominio con fondo irregular, transformamos el espacio físico $(x, z)$ al espacio rectangular computacional $(x, \eta)$ mediante:
\begin{equation}
\eta = \frac{z}{h(x)} \quad \implies \quad \eta \in [0, 1]
\end{equation}
donde $\eta = 0$ representa la base ($z=0$) y \(\eta = 1\) corresponde a la superficie libre de la duna ($z=h(x)$).

Las derivadas parciales se transforman aplicando la regla de la cadena:
\begin{align}
\frac{\partial}{\partial z} &= \frac{1}{h(x)} \frac{\partial}{\partial \eta} \\
\frac{\partial}{\partial x}\Big|_z &= \frac{\partial}{\partial x}\Big|_{\eta} - \frac{\eta}{h(x)} \frac{dh}{dx} \frac{\partial}{\partial \eta}
\end{align}

---

\section{Campo de Velocidades Modulado y Libre de Divergencia}
Para generar la deposición periódica por capas en el sotavento, modulamos la función de corriente del flujo principal $\psi_{\text{bulk}}(x, \eta, t)$ de manera temporal y espacial.

\subsection{Definición de la Función de Corriente (\(\psi\))}
Definimos la función de corriente en coordenadas sigma como:
\begin{equation}
\psi(x, \eta, t) = U_0 H_{\text{base}} \eta^{m+1} g(x, t) - c_{\text{mig}} h(x) \eta
\end{equation}
donde $U_0$ es la velocidad de referencia, $m$ es el exponente de ley de potencia del perfil de velocidades y $g(x, t)$ es la función de modulación temporal y espacial:
\begin{equation}
g(x, t) = 1.0 + A \sin\left(\frac{2\pi t}{T}\right) S(x)
\end{equation}
Aquí, $A$ es la amplitud de la fluctuación del flujo, $T$ es el periodo de deposición de las capas y $S(x)$ es una función de transición suave (sigmoide) que activa la oscilación únicamente en sotavento:
\begin{equation}
S(x) = \frac{1}{2} + \frac{1}{2}\tanh\left(\frac{x - x_{\text{crest}}}{\delta}\right)
\end{equation}
con un ancho de transición $\delta = 0.02$. La derivada espacial del sigmoide es:
\begin{equation}
\frac{dS}{dx} = \frac{1}{2\delta} \text{sech}^2\left(\frac{x - x_{\text{crest}}}{\delta}\right)
\end{equation}

\subsection{Velocidades Derivadas}
Las velocidades en coordenadas sigma se deducen a partir de la función de corriente para garantizar la conservación de la masa del fluido:

\begin{enumerate}
    \item \textbf{Velocidad Horizontal (\(u\))}:
    \begin{equation}
    u(x, \eta, t) = \frac{1}{h(x)} \frac{\partial \psi}{\partial \eta} = \frac{U_0 H_{\text{base}}}{h(x)} (m+1) \eta^m \left[ 1.0 + A \sin\left(\frac{2\pi t}{T}\right) S(x) \right] - c_{\text{mig}}
    \end{equation}
    
    \item \textbf{Velocidad Vertical de Coordenadas (\(\omega\))}:
    \begin{equation}
    \omega(x, \eta, t) = -\frac{\partial \psi}{\partial x} = c_{\text{mig}} \eta \frac{dh}{dx} - U_0 H_{\text{base}} \eta^{m+1} A \sin\left(\frac{2\pi t}{T}\right) \frac{dS}{dx}
    \end{equation}
\end{enumerate}

\subsection{Verificación de la Divergencia Cero}
La divergencia en coordenadas sigma está dada por:
\begin{equation}
\nabla \cdot \vec{u} = \frac{\partial (h u)}{\partial x} + \frac{\partial \omega}{\partial \eta}
\end{equation}
Derivando las expresiones de $h u$ y $\omega$:
\begin{align}
\frac{\partial (h u)}{\partial x} &= U_0 H_{\text{base}} (m+1) \eta^m A \sin\left(\frac{2\pi t}{T}\right) \frac{dS}{dx} - c_{\text{mig}} \frac{dh}{dx} \\
\frac{\partial \omega}{\partial \eta} &= c_{\text{mig}} \frac{dh}{dx} - U_0 H_{\text{base}} (m+1) \eta^m A \sin\left(\frac{2\pi t}{T}\right) \frac{dS}{dx}
\end{align}
Sumando ambos términos se obtiene:
\begin{equation}
\nabla \cdot \vec{u} = 0
\end{equation}
Esto demuestra que el campo de velocidades es exactamente libre de divergencia para todo tiempo $t$, conservando de forma estricta la masa de la mezcla.

---

\section{Ecuación de Advección-Segregación}
La conservación de la fracción de partículas finas $\phi_s(x, \eta, t)$ está gobernada por:
\begin{equation}
\frac{\partial Q}{\partial t} = -\frac{\partial (u Q)}{\partial x} - \frac{\partial}{\partial \eta} \left( \omega \phi_s + F_{\text{seg}} \right)
\end{equation}
donde:
\begin{itemize}
    \item $Q = h(x) \phi_s$ es la cantidad conservada.
    \item $u Q$ es el flujo advectivo horizontal.
    \item $\omega \phi_s$ es el flujo advectivo vertical en coordenadas sigma.
    \item $F_{\text{seg}} = - q_{\text{seg}} \phi_s (1 - \phi_s)$ es el flujo de segregación por gravedad ($q_{\text{seg}}$ es la velocidad máxima de segregación).
\end{itemize}

---

\section{Condiciones de Borde Dinámicas}
La conservación de la masa total se controla mediante el balance entre la erosión en barlovento y la deposición en sotavento en la superficie libre ($\eta = 1$).

\subsection{Erosión en Barlovento (\(\omega \ge 0\))}
Flujo saliente:
\begin{equation}
F_{\eta}(x, 1) = \omega \phi_s(x, 1)
\end{equation}

\subsection{Deposición en Sotavento (\(\omega < 0\))}
Flujo entrante. La concentración de alimentación $\phi_{\text{inflow}}(x, t)$ se calcula dinámicamente como:
\begin{equation}
\phi_{\text{inflow}}(x, t) = \text{clip}(\lambda(t) P(x, t), 0.0, 1.0)
\end{equation}

Para simular la alternancia granulométrica, la función de peso $P(x, t)$ se modula doblemente (espacial y temporalmente):
\begin{equation}
P(x, t) = P_{\text{spatial}}(x) \cdot \left[ 1.0 + A_P \sin\left(\frac{2\pi t}{T}\right) \right]
\end{equation}
donde $A_P$ es la amplitud de la oscilación de la clasificación, y $P_{\text{spatial}}(x)$ es:
\begin{equation}
P_{\text{spatial}}(x) = \max\left(0.1, 1.0 + \alpha_{\text{spatial}} \left( 0.5 - \frac{x - x_{\text{crest}}}{L - x_{\text{crest}}} \right)\right)
\end{equation}

El factor de escala global $\lambda(t)$ se calcula resolviendo el balance de flujos netos en la superficie:
\begin{equation}
\lambda(t) = \frac{\int_{\text{erosión}} \omega(x) \phi_s(x, 1) \, dx}{\int_{\text{deposición}} |\omega(x)| P(x, t) \, dx}
\end{equation}

---

\section{Esquema de Integración Temporal (SSP-RK2)}
Para resolver la ecuación de advección con campos dependientes del tiempo, el paso de tiempo se realiza mediante un esquema Runge-Kutta de dos pasos que preserva la estabilidad fuerte (SSP-RK2):

\begin{align}
Q^1 &= Q^n + \Delta t \, \text{RHS}(Q^n, t^n) \\
Q^{n+1} &= \frac{1}{2} Q^n + \frac{1}{2} \left( Q^1 + \Delta t \, \text{RHS}(Q^1, t^n + \Delta t) \right)
\end{align}
donde $\text{RHS}$ representa las derivadas espaciales (advección upwind y segregación). Posterior a cada paso RK2, se aplica una corrección de masa global para garantizar de manera estricta que $\sum Q \Delta x \Delta \eta$ se conserve a nivel de precisión de máquina.

\end{document}
